% ============================================================
% 第 6 章  电力系统三相短路电流的实用计算
% （教材：何仰赞、温增银《电力系统分析》上册 第 6 章）
% 结构对齐 6.1—6.4；例题、习题不录
% ============================================================

\chapter{电力系统三相短路电流的实用计算}

第 5 章解决了“短路电流随时间怎样变化”，本章解决“在给定系统中怎样算出来”。
实用计算有两条腿：\textbf{节点阻抗矩阵法}（计算机路线，精确、通用）与
\textbf{转移阻抗 $+$ 计算曲线法}（手算路线，查表、快速）。
两者的共同前提是：用次暂态电势 $E''$ 与次暂态电抗 $X_d''$ 作发电机的等值电路，
并把网络化为纯电抗网络（略电阻、略线路电容、略变压器励磁支路）。

\section{短路电流计算的基本原理和方法}

\subsection{电力系统节点方程的建立}

短路计算关心的是\textbf{起始次暂态电流}，因此各发电机用 $E''$、$X_d''$ 表示，
负荷一般用恒定阻抗（或按“略去”处理）。选定基准值后作等值网络，
按第 4 章的方法组装节点导纳矩阵 $\boldsymbol Y$，得节点方程
\begin{equation}
  \boldsymbol Y\dot{\boldsymbol U}=\dot{\boldsymbol I},
  \qquad \dot I_i=\frac{\dot E_i''}{jX_{di}''}\quad(\text{电源节点}),
  \qquad \dot I_i=0\quad(\text{无源节点}).
\end{equation}
求解前先在发电机节点“注入”各电源的短路电流，得到\textbf{短路前}的节点电压
$\dot U_i(0)$（也可直接取 $\dot U_i(0)\approx1\angle0^\circ$）。

\begin{table}[htbp]
  \centering
  \caption{短路计算中各类元件的处理}
  \small
  \begin{tabular}{@{}p{0.18\linewidth}p{0.34\linewidth}p{0.36\linewidth}@{}}
    \toprule
    元件 & 处理方式 & 理由 \\
    \midrule
    发电机 & 电势 $\dot E''$ 串 $jX_d''$ & 短路瞬间阻尼绕组磁链不变（第 5 章） \\
    变压器 & 只有电抗，励磁支路略去 & 励磁电流小 \\
    线路 & 只有电抗，并联电纳略去 & 短路计算中容性电流次要 \\
    负荷 & 恒定阻抗，或按“略去”处理 & 负荷电流远小于短路电流 \\
    无限大电源 & 内电抗为零，端电压恒定 & 系统容量远大于所研究部分 \\
    \bottomrule
  \end{tabular}
\end{table}

\subsection{利用节点阻抗矩阵计算短路电流}

设在节点 $f$ 发生短路，用一条阻值为 $z_f$ 的故障支路接地（$z_f=0$ 为金属性短路）。
短路后节点 $f$ 的电压为该支路两端电压，短路电流即流过 $z_f$ 的电流。
由叠加原理（把故障看作在 $f$ 点叠加一个 $-U_f(0)/z_f$ 的注入电流），可得
\keyformula{\dot I_f=-\frac{\dot U_f(0)}{Z_{ff}+z_f}}
金属性短路（$z_f=0$）时 $\dot I_f=-\dot U_f(0)/Z_{ff}$，其模为 $U_f(0)/Z_{ff}$。
各节点电压随之为
\begin{equation}
  \dot U_i=\dot U_i(0)-Z_{if}\dot I_f
  \qquad(i=1,2,\dots,n).
\end{equation}
仅需 $Z_{ff}$（输入阻抗）与一列 $Z_{if}$（转移阻抗），不必求全逆——这正是第 4.3 节
“只求一列”的动机。

若短路点不是原网络节点，可在该处增设一个节点（追加树支），再按上式算。

\subsection{利用电势源对短路点的转移阻抗计算短路电流}

把网络化简为“以短路点 $f$ 为中心、各电源点为顶点的星形电路”，则第 $i$ 个电源
对 $f$ 的支路电抗 $X_{if}$ 称\textbf{转移阻抗}，其电流为
\keyformula{\dot I_i=\frac{\dot E_i}{Z_{if}},\qquad \dot I_f=\sum_{i=1}^{m}\dot I_i}
\textbf{转移阻抗的物理意义}：除第 $i$ 个电源外其余电源电势为零时，
$f$ 点短路电流与 $E_i$ 之比；等价地，$Z_{if}$ 是“$i$ 点注入单位电流时 $f$ 点的电压”。

转移阻抗与输入阻抗的关系（同一网络、全电源同电势时两种算法结果必须相等）：
\begin{equation}
  \frac{1}{Z_{ff}}=\sum_{i=1}^{m}\frac{1}{Z_{if}} .
\end{equation}
这一恒等式是校验网络化简是否正确的常用手段。

\begin{table}[htbp]
  \centering
  \caption{两种算路方法的比较}
  \small
  \begin{tabular}{@{}p{0.16\linewidth}p{0.32\linewidth}p{0.34\linewidth}@{}}
    \toprule
    比较项 & 节点阻抗矩阵法 & 转移阻抗法 \\
    \midrule
    适用 & 复杂网络、多点短路、需各节点电压 & 手算、电源少、需分电源电流 \\
    关键量 & $Z_{ff}$ 与一列 $Z_{if}$ & 各 $Z_{if}$（星形化） \\
    优点 & 一次求逆可算所有节点短路 & 物理图像清晰、便于查计算曲线 \\
    缺点 & 需计算机 & 网络大时化简繁琐 \\
    \bottomrule
  \end{tabular}
\end{table}

\corepoint{“转移阻抗”与“输入阻抗”不是一回事：$Z_{ff}$ 是 $f$ 点看进去的等值阻抗，
$Z_{if}$ 是第 $i$ 个电源到 $f$ 的等值阻抗。二者只在单电源网络中数值相等。
多电源时若误用 $Z_{ff}$ 去算某一台机的电流，就会把其他电源的贡献也算进去。}

\section{起始次暂态电流和冲击电流的实用计算}

\subsection{起始次暂态电流}

实用计算中取各发电机 $\dot E''\approx1.0\sim1.1$（标幺），工程上常直接取 $1$，
于是短路点电流的标幺值为
\keyformula{I''_{*}=\frac{1}{X_{\Sigma *}},\qquad I''=I''_{*}I_B=\frac{I''_{*}S_B}{\sqrt{3}U_{av}}}
$X_{\Sigma*}$ 为从短路点看进去的等值电抗（标幺），$U_{av}$ 取短路点所在电压级的平均额定电压。

\textbf{关于负荷的处理}：负荷在短路瞬间可能“变成电源”向短路点供电（异步电动机的反馈电流），
判断依据是\textbf{残余电压}：
\begin{itemize}
  \item 若短路后该负荷节点残余电压 $U>0.8U_N$，负荷不供给短路电流，可略去；
  \item 若 $U<0.8U_N$，异步电动机可能倒送电流，应按 $E''\approx0.8$、$X''\approx0.2$ 的
        电动机模型计入（其冲击系数取 $1$）。
\end{itemize}

\subsection{冲击电流的实用计算}

\begin{equation}
  i_{sh}=K_{sh}\sqrt{2}I'' ,
\end{equation}
冲击系数按短路点位置与电源类型选取：

\begin{table}[htbp]
  \centering
  \caption{冲击系数的取法}
  \small
  \begin{tabular}{@{}lll@{}}
    \toprule
    短路点位置 & $K_{sh}$ & $i_{sh}$ \\
    \midrule
    发电机出口或发电机电压母线 & $1.90$ & $2.69I''$ \\
    高压电网（远离发电机） & $1.80$ & $2.55I''$ \\
    由异步电动机供给的分量 & $1.00$ & $\sqrt{2}I''$ \\
    \bottomrule
  \end{tabular}
\end{table}

继电保护整定和热稳定校验用 $I''$；动稳定校验用 $i_{sh}$ 与 $I_{sh}=1.52I''$（或 $1.62I''$）。

\section{短路电流计算曲线及其应用}

\subsection{计算曲线的概念}

由第 5 章，短路电流周期分量是许多参数的复杂函数：各发电机的电抗与时间常数、
反映短路前运行状态的电势初值、励磁系统参数、短路点离机端的距离、时间 $t$。
在发电机（含励磁系统）参数和运行初态给定后，它只是“电气距离”和“时间”的函数。
把\textbf{归算到发电机额定容量}的外接电抗标幺值与发电机纵轴次暂态电抗标幺值之和定义为
\textbf{计算电抗}
\keyformula{X_{js}=X_d''+X_e}
则短路电流周期分量的标幺值可写成
\begin{equation}
  I_{p*}=f\left(X_{js},t\right),
\end{equation}
反映这一函数关系的一族曲线（或数字表）即\textbf{计算曲线}。

\subsection{计算曲线的制作条件}

计算曲线不是严格的理论结果，而是在一组\textbf{典型条件}下对多台样机逐台计算再取平均得到的。
用曲线前必须知道它是在什么条件下做出来的，否则“查表”就失去依据：

\begin{table}[htbp]
  \centering
  \caption{计算曲线的制作条件}
  \small
  \begin{tabular}{@{}p{0.20\linewidth}p{0.68\linewidth}@{}}
    \toprule
    项目 & 条件 \\
    \midrule
    典型接线 & 发电机经变压器接到高压母线，$50\%$ 负荷接于发电厂高压母线，
               其余负荷经线路送到短路点以外 \\
    短路前状态 & 发电机额定满载运行 \\
    负荷模型 & 恒定阻抗 $Z_{LD}=U^{2}/\left(S_{LD}(\cos\varphi+j\sin\varphi)\right)$，
               取 $U=1$、$\cos\varphi=0.9$ \\
    励磁 & 均配强行励磁，顶值电压为额定励磁电压的 $1.8$ 倍；
           励磁系统等值时间常数 $T_e$：汽轮发电机 $0.25$~s，水轮发电机 $0.02$~s \\
    样机 & 汽轮发电机取 $12\sim200$~MW 的 18 种型号；水轮发电机取 $12.5\sim225$~MW 的 17 种型号；
           对给定 $X_{js}$ 与 $t$ 取各机计算结果的算术平均值 \\
    适用范围 & 只做到 $X_{js}=3.45$；$X_{js}\ge3.45$ 时认为周期分量幅值不随时间变化，
               直接取 $I_{p*}=1/X_{js}$ \\
    \bottomrule
  \end{tabular}
\end{table}

制作曲线的网络只含\textbf{一台}发电机，且计算电抗与负荷支路无关。因此实际系统用曲线前，
必须先把网络化为“以短路点为中心、以各电源点为顶点的星形电路”。

\subsection{计算曲线的应用}

\paragraph{（1）绘制等值网络}
取 $S_B$、$U_B=U_{av}$；发电机电抗用 $X_d''$；略去各元件电阻、线路电容和变压器励磁支路；
无限大电源内电抗取零；略去负荷。

\paragraph{（2）网络变换、合并电源}
按“短路电流变化规律相近”的原则合并：与短路点电气距离相近的\textbf{同类型}发电机可合并；
远离短路点的同类型电厂可合并；\textbf{直接接于短路点的发电机应单独考虑}；
无限大电源必须单独计算（其电流不衰减）。合并后得 $g$ 组等值机，求出各组的转移电抗
$X_{if}$ 与无限大电源的 $X_{fs}$。

\paragraph{（3）归算计算电抗}
\begin{equation}
  X_{jsi}=X_{if}\frac{S_{Ni}}{S_B}\qquad(i=1,2,\dots,g),
\end{equation}
$S_{Ni}$ 为第 $i$ 台等值机所代表的那部分发电机额定容量之和。

\paragraph{（4）查曲线}
由 $X_{jsi}$ 查相应类型（汽轮/水轮）的计算曲线，得指定时刻 $t$ 的标幺值 $I_{pti*}$。

\paragraph{（5）无限大电源的分量}
\begin{equation}
  I_{ps*}=\frac{1}{X_{fs}} .
\end{equation}

\paragraph{（6）换算为有名值}
\begin{equation}
  I_{ti}=I_{pti*}\frac{S_{Ni}}{\sqrt{3}U_{av}},\qquad
  I_{s}=I_{ps*}\frac{S_B}{\sqrt{3}U_{av}},\qquad
  I_t=\sum_{i=1}^{g}I_{ti}+I_s ,
\end{equation}
其中 $U_{av}$ 取短路处电压级的平均额定电压，$S_{Ni}/(\sqrt{3}U_{av})$ 即归算到短路处电压级的
第 $i$ 台等值机额定电流。

\begin{table}[htbp]
  \centering
  \caption{手算路线的三个“归算”易混点}
  \small
  \begin{tabular}{@{}p{0.24\linewidth}p{0.30\linewidth}p{0.38\linewidth}@{}}
    \toprule
    量 & 归算到什么容量 & 用途 \\
    \midrule
    转移电抗 $X_{if}$ & 统一基准 $S_B$ & 描述网络拓扑，用于合并化简 \\
    计算电抗 $X_{jsi}$ & 各电源自己的 $S_{Ni}$ & 查计算曲线（曲线的横坐标） \\
    电流标幺值 $I_{pti*}$ & 各电源自己的 $S_{Ni}$ & 乘 $S_{Ni}/(\sqrt{3}U_{av})$ 得有名字 \\
    \bottomrule
  \end{tabular}
\end{table}

\corepoint{计算曲线的横坐标是“按本机容量归算的”计算电抗，而网络化简得到的转移电抗是
“按 $S_B$ 归算的”。二者差一个 $S_{Ni}/S_B$ 的因子，漏掉这一步是最常见的错误。
另外 $X_{js}\ge3.45$ 时曲线终止，此时直接取 $1/X_{js}$，不要再外推查表。}

\section{短路电流周期分量的近似计算}

工程估算中可假定短路电路接于\textbf{内阻抗为零的恒电势电源}，于是周期分量幅值不随时间变化，
只有非周期分量衰减。取各电源电势标幺值为 $1$，则
\keyformula{I_{p*}=\frac{1}{X_{ff*}},\qquad I_p=I_{p*}I_B=\frac{I_{p*}S_B}{\sqrt{3}U_{av}},\qquad
  S_k=\frac{S_B}{X_{ff*}}}
$X_{ff*}$ 为短路点的输入电抗（标幺）。这样算出的短路电流（或短路功率）比实际偏大，
但差别随短路点距离增大而迅速减小——因为短路点越远，“电源电压恒定”的假设越接近实际，
尤其当发电机装有自动励磁调节器时更是如此。故这一简化可用来估计短路电流
（或短路功率）的\textbf{最大可能值}。

\paragraph{未知系统参数的估算}
若只知某母线处的短路容量 $S_s$（或有名值电流 $I_s$），可把该系统等效为经电抗 $X_s$
接于该母线的无限大电源：
\begin{equation}
  X_{s*}=\frac{S_B}{S_s}\qquad\left(\text{或由}\ S_s=\sqrt{3}U_{av}I_s\ \text{反算}\right).
\end{equation}
若连短路容量也不知道，可由该处变电所断路器的\textbf{额定切断容量}反推：
认为断路器后三相短路时其切断容量恰好被充分利用，即 $S_s$ 取断路器额定切断容量。

\begin{table}[htbp]
  \centering
  \caption{本章四种算法的适用场合}
  \small
  \begin{tabular}{@{}p{0.22\linewidth}p{0.30\linewidth}p{0.36\linewidth}@{}}
    \toprule
    方法 & 精度 & 适用 \\
    \midrule
    节点阻抗矩阵法 & 高（取决于 $E''$ 取值） & 计算机计算、需全网点电压 \\
    转移阻抗 $+$ 计算曲线 & 中（受制作条件限制） & 手算、需分电源、需任意时刻电流 \\
    近似计算（$E=1$） & 偏低精度、偏保守 & 估算最大短路电流、初选设备 \\
    按短路容量反推系统电抗 & 粗 & 缺少系统详细数据时 \\
    \bottomrule
  \end{tabular}
\end{table}

\corepoint{本章的核心是一条计算链：
\textbf{选基准 $\to$ 作纯电抗等值网络 $\to$ 化简到短路点 $\to$ 求 $X_{\Sigma}$（或 $X_{js}$）
$\to$ 得 $I''$ $\to$ 派生 $i_{sh}$、$I_{sh}$、$S_k$}。
链条上每一环都只做一件事；出错往往不在物理，而在“归算到哪个容量”“用哪一电压级的
$U_{av}$”这两个纯约定问题上。}
